\chapter{Mise en oeuvre, validation et perspectives}

\section{Démarche de mise en oeuvre}

La mise en oeuvre d'un nouveau modèle d'alerting en production impose une approche incrémentale.

La démarche retenue s'articule autour de trois étapes :
\begin{itemize}
	\item qualifier l'existant et prioriser les alarmes à forte valeur,
	\item introduire progressivement les nouvelles règles de criticité,
	\item observer les effets en exploitation avant généralisation.
\end{itemize}

Cette méthode limite les risques de régression et facilite l'appropriation par les équipes d'astreinte.

\section{Stratégie de validation}

La validation d'une évolution d'alerting ne se limite pas à la réussite technique du déploiement. Elle doit mesurer
la pertinence opérationnelle du système obtenu.

Les indicateurs à suivre sont notamment :
\begin{itemize}
	\item volume d'alertes par niveau de criticité,
	\item ratio alertes actionnables / alertes non pertinentes,
	\item temps moyen de détection et d'orientation,
	\item ressenti des équipes d'astreinte (fatigue, confiance dans les alertes).
\end{itemize}

Ces métriques permettent d'ajuster les seuils et la classification dans une logique d'amélioration continue.

\section{Limites et perspectives}

Le travail réalisé constitue une base méthodologique et technique solide, mais plusieurs axes restent ouverts :
\begin{itemize}
	\item affiner la granularité des niveaux de criticité selon les services,
	\item renforcer l'intégration avec les runbooks d'incident,
	\item améliorer encore le ciblage des notifications par rôle d'astreinte,
	\item industrialiser le suivi de qualité des alertes dans la durée.
\end{itemize}

Ces perspectives visent un système d'alerting plus robuste, plus lisible et mieux aligné avec les besoins réels des
équipes d'exploitation.